\section{Power fusion in the alternative tiling certificate}
\label{sec:power-fused-tiling}

This section gives an arithmetic improvement, not a new numerical bound for
ordinary Diophantine equations. For every fixed finite strict-boundary Wang
tile set, the previous five-power interface can be replaced by either
seven exposed natural coordinates and three power graphs, or eight exposed
natural coordinates and two power graphs. All graph witnesses remain to be
compiled and counted. The construction does not require Jones's universal
formula, K215, or restriction to square boards.

\subsection{Recovering geometry instead of guessing its exponents}
Fix $k\geq1$ tiles with edge colors $(w_t,e_t,n_t,s_t)$. Color zero is
required exactly on the outer boundary. Choose a power of two $C=2^c$
greater than $k$ and all colors, put $A_0=4C^2$, and choose a power of two
$D\geq32kA_0$. Initially use natural coordinates $a,Q,G,Z,r$ and the graph
\[
 P=1+(C-1)G=2^{c(a+1)}.
\]
Thus $P\geq C$ and $G\geq1$. Set
\[
 H=A_0P,\qquad R=DP^3,\qquad B=R/2,\qquad L=R^{2k}.
\]
For a formal vector $x=(x_0,\ldots,x_{k-1})$, write $U=\sum_t x_t$;
$W_0,E_0,N_0,S_0$ denote its sums over tiles having the indicated edge
color zero. The seven residuals $\ell_j(x)=\sum_t a_{jt}x_t+c_j$ are
\[
\begin{aligned}
&U-G,\quad (Q-1)W_0-(P-1),\quad QW_0-CE_0,\\
&C\sum_t e_tx_t-\sum_t w_tx_t,\quad
 Q\sum_t s_tx_t-\sum_t n_tx_t,\\
&Q-1-(C-1)N_0,\quad QS_0-PN_0.
\end{aligned}
\]
Here $x$ is not a tuple of certificate variables. It is recovered from $Z$.
The third and sixth residuals have been reversed relative to the predecessor
so that unbounded positive $Q$ cannot create negative packed summands.
Define ordinary polynomial expressions
\[
\begin{aligned}
 K_j&=\sum_{t=0}^{k-1}(H+a_{jt})R^{k-1-t},\\
 F_j&=K_jZ+(B-HG+c_j)R^{k-1},\\
 T_7&=F_7+Q\frac{L}{2P},\\
 M_0&=G\sum_{t=0}^{k-1}R^t,\\
 M_1&=R^{2k-1}-1-(R-1-B)R^{k-1},\\
 M_*&=L-1-(R-1-B)R^{k-1},\\
 A&=M_0+M_1\sum_{j=1}^{6}L^j+M_*L^7,\\
 Y&=Z+\sum_{j=1}^{6}F_jL^j+T_7L^7.
\end{aligned}
\]
The apparent quotient $L/(2P)$ means the ordinary polynomial
$(D^{2k}/2)P^{6k-1}$. Impose the total-binomial graph
$\binom AY=2r+1$, with value zero when $Y>A$.

\paragraph{Bootstrap.}
Before bounding $Q$, every $H+a_{jt}$ is positive, every offset is at
least $B-HP-P>0$, and $F_7\geq Z$. Binary Lucas parity gives
$Y\subseteq_{\rm bits}A<L^8$. Therefore
\[
 T_7<L,\qquad Z<L,\qquad Q<2P.
\]
The bottom $L$-block now gives $Z\subseteq_{\rm bits}M_0$, so
$Z=\sum_{t<k}x_tR^t$ with $x_t\subseteq_{\rm bits}G$ and $x_t<P$.
This establishes the bounds before using any no-carry argument.

After these bounds, $|a_{jt}|\leq H/2$ and $|c_j|\leq2P$.
Noncentral convolution coefficients are less than $2kHG<R/8$.
The central coefficient is
\[
 d_j=B+\ell_j(x)+H(U-G),\qquad
 0<B-HG-2P\leq d_j\leq B+2kHG+2P<R.
\]
Thus $F_j<R^{2k-1}$ with no carries. The guard occupies only the unused
highest digit, because
$QL/(2P)=(QR/(2P))R^{2k-1}$ and $QR/(2P)<R$.
All lanes separate. Their central masks allow only the bit $B$, so
$d_j=B$. The first row gives $(H+1)(U-G)=0$, and then every $\ell_j=0$.

\paragraph{Geometry.}
Write $P=C^N$. The mask $G=1+C+\cdots+C^{N-1}$ and $k<C$ make
$U=G$ an exact one-tile-per-cell identity. Boundary rows imply
\[
 (Q-1)W_0=P-1>0,\quad (C-1)N_0=Q-1,\quad QS_0=PN_0.
\]
Hence $Q\geq2$, $\gcd(Q,N_0)=1$, and $Q\mid P$. Since $P$ is a
power of two, $Q=2^u$; $C-1\mid Q-1$ gives $c\mid u$, so $Q=C^w$.
Also $Q-1\mid P-1$ gives $w\mid N$. These assertions use the elementary
identity $b^s-1\mid b^t-1$ if and only if $s\mid t$ for positive $s,t$.
Taking $h=N/w$ recovers the rectangle, including its exact four boundary
masks. The two color residuals check all adjacent edges. Conversely, a
rectangle gives these masks and $a=wh-1$; its guard is at most $L/2$,
while $F_7<R^{2k-1}$, so it satisfies the packed condition.
This proves the five-coordinate, one-power/one-binomial interface for
arbitrary rectangles.

\subsection{A three-power interface with seven coordinates}
Replace $G$ by $g+1$, so $P=C+(C-1)g\geq C$, and keep the same circuit.
Use the seven coordinates $Q,g,Z,v,q,s,t$ and require
\[
\begin{aligned}
 Pv&=2^{A+1},\\
 s+t+1&=2^{(A+1)Y},\\
 (s+t+1)(2q+1)+s&=(Pv+1)^A.
\end{aligned}
\tag{T3}
\]
The first graph makes $P$ a power of two and therefore a power of $C$.
For $d=2^{A+1}$, all binomial coefficients of $(d+1)^A$ are smaller
than $d$. The last two equations say that its quotient by $d^Y$ is odd,
with a legitimate remainder $s<d^Y$. Thus they are exactly the binomial
parity condition, including rejection of $Y>A$.

For completeness, a genuine tiling has $Y\leq A$, $P\leq A$, and
$A\geq10$. These bounds follow directly from the highest digit of $M_*$
and $R\geq512P^3$. Thus $2^{A+1}/P$ is an integer; choose it as $v$
and use the Euclidean quotient and remainder. Geometry's separate power
graph and area index have both disappeared.

\subsection{Exact synchronization: two graphs and eight coordinates}
The following elementary lemma removes another power graph. Suppose
$n,Y,d,e$ are positive, $d>n+1$, $1\leq Y\leq n$, $e=2^n$, $d\mid e$,
and
\[
 e-1=(d-1)\bigl(Y+(d-1)j\bigr).
\]
Then $d=2^u$ and $d-1\mid e-1$ give $e=d^h$ with $h\leq n$.
The displayed equation gives $h\equiv Y\pmod{d-1}$, so the bound forces
$h=Y$. Without the size bound the conclusion is false: $e=8,d=2,Y=1$
satisfies the congruence with $j=6$.

Take the eight natural coordinates $Q,g,Z,v,q,u,t,j$ and define
\[
\begin{aligned}
 n&=(A+1)Y,& d&=P(n+v+1),\\
 e&=d(u+t+1),& W&=e(2q+1)+du+1,\\
 S&=e-1-(d-1)\bigl(Y+(d-1)j\bigr).
\end{aligned}
\]
Require just the two graphs and one ordinary residual
\[
 e=2^n,\qquad W=(d+1)^A,\qquad S=0.
\tag{T2}
\]
Pre-bootstrap positivity gives $Y\geq1$; $P\geq2$ gives
$d\geq2(n+1)>n+1$. The lemma yields $e=d^Y$, whence
$d=2^{A+1}$. Since $P\mid d$, the required geometry follows. Moreover
$0\leq du+1<e$, as $e-(du+1)=d(t+1)-1\geq1$. The second graph
therefore has odd Euclidean quotient, exactly the binomial parity test.

For the converse set $d=2^{A+1}$, $e=d^Y$, and write
$(d+1)^A=em+s$ with $0\leq s<e$. Here $m$ is positive odd and
$s\equiv1\pmod d$. Choose
\[
 q=\frac{m-1}{2},\quad u=\frac{s-1}{d},\quad
 t=\frac ed-u-1,\quad
 j=\frac{(e-1)/(d-1)-Y}{d-1},\quad v=\frac dP-n-1.
\]
All are nonnegative integers. For the only delicate bound, use
$Y\leq A$, $P\leq A$, $A\geq10$, and
\[
 2^{A+1}>A^3+A^2+A\geq P\bigl((A+1)Y+1\bigr).
\]
The exponential inequality holds at $A=10$ and propagates by induction.
This proves the eight-coordinate, two-power interface.

\paragraph{Honest ordinary-polynomial accounting.}
If exact fixed-base and general-base graph polynomials have respectively
$p_2,p_g$ internal natural witnesses, and total binomial has $b$, the
three interfaces have respectively
\[
 5+p_2+b,\qquad 7+2p_2+p_g,\qquad 8+p_2+p_g
\]
ordinary-polynomial variables after sum-of-squares combination. In the last
case the literal template is $E_2(n,e)^2+E_g(d+1,A,W)^2+S^2=0$, with
independent internal witness tuples. Since $\deg A,\deg Y\leq48k$, its
degree is at most
\[
 2\max\{(96k+2)\delta_2,(96k+3)\delta_g,192k+3\}.
\]
The missing graph implementations, direct ZFC-checker tile instance and
internal ZFC interpretation remain explicit obligations. The supplied
standard-library checks and fixed-power expression DAGs validate the
arithmetic implementation but are not a proof-assistant formalization.
